\documentclass[11pt]{article}

% Change "review" to "final" to generate the final (sometimes called camera-ready) version.
% Change to "preprint" to generate a non-anonymous version with page numbers.
\usepackage[review]{acl}

\usepackage{times}
\usepackage{latexsym}
\usepackage[T1]{fontenc}
\usepackage[utf8]{inputenc}
\usepackage{microtype}
\usepackage{graphicx}
\usepackage{amsmath}
\usepackage{amssymb}
\usepackage{booktabs}
\usepackage{multirow}
\usepackage{array}
\usepackage{xcolor}

% ---- macros -----------------------------------------------------------
\newcommand{\Htraj}{\mathcal{H}}
\newcommand{\Htogo}{H_{\mathrm{togo}}}
\newcommand{\AH}{A_{H}}
\newcommand{\Om}{\Omega}
\newcommand{\Omt}{\tilde{\Omega}}
\newcommand{\E}{\mathbb{E}}
\newcommand{\gradt}{\nabla_{\theta}}
\newcommand{\pit}{\pi_{\theta}}
\newcommand{\TODO}[1]{\textbf{[TODO: #1]}}

\title{Entropy Collapse Has a Direction:\\
Correcting the Visitation Channel in Token-Level Entropy Control for RLVR}

\author{Anonymous ACL submission}

\begin{document}
\maketitle

\begin{abstract}
Reinforcement learning with verifiable rewards (RLVR) reliably drives a policy's
entropy toward zero, and a growing family of methods intervenes by re-weighting the
policy loss with a per-token estimate of how much each token's update costs in
entropy. We show that every such method estimates only half of the relevant
quantity. The gradient of the \emph{trajectory} entropy splits exactly into two
terms: a \emph{local} channel, which asks how the conditional distribution at a
visited state flattens or sharpens, and a \emph{visitation} channel, which asks how
the probability of \emph{reaching} that state changes. Existing token-level
estimators capture the first and are structurally blind to the second. We show the
blindness is worse than an omission: the local estimator is exactly zero at a
uniform branch point, that zero is the correct first-order answer, and the
visitation term's prefactor is simultaneously maximal there --- the two channels are
\emph{anti-correlated}, so no refinement of the local term can recover the missing
signal. We then derive an estimator for the visitation channel that costs one extra
forward pass. The reward-to-go implied by the decomposition is the remaining
conditional entropy $\Htogo$, which we forecast with parallel
multi-token-prediction heads; the sibling-prefix baseline that keeps the estimator
unbiased is supplied by a tree-structured rollout in which several samples share a
prefix and diverge at scheduled depths. The resulting advantage $\AH$ sums to zero
within a sibling set by construction, so the method cannot change the average amount
of damping --- it can only decide which of two siblings is damped. We give the full
derivation, the four approximations it requires and what each discards, and a
five-arm design that isolates the contribution of the forecast's \emph{value} from
that of the apparatus carrying it.
\end{abstract}

\section{Introduction}
\label{sec:intro}

Policy-gradient fine-tuning with verifiable rewards sharpens a language model's
output distribution. Measured as the mean per-token entropy of the policy, this
sharpening is fast, monotone, and, past a point, coincident with the end of useful
learning. The standard reading is that the collapse is the problem, and the standard
remedy is to spend part of the update budget slowing it down: an entropy bonus, a KL
anchor, a clip on the tokens that lose the most entropy, or --- in the most recent
line of work --- a per-token re-weighting of the policy loss by a first-order
estimate of the entropy each token's update destroys.

We argue that the accounting behind that remedy is incomplete, and that the
incompleteness is structural rather than a matter of estimator quality.

Write $\Htraj(\pit \mid x)$ for the entropy of the distribution the policy induces
over whole responses to a prompt $x$. The autoregressive factorization turns it into
an expectation of a sum of per-token conditional entropies, and the gradient of that
expression splits, by the product rule, into exactly two terms. The first holds the
visited states fixed and asks how the conditional distribution at each one changes:
this is what a per-token entropy estimate measures. The second holds the per-state
entropies fixed and asks how the probability of \emph{visiting} each state changes.
An update that concentrates mass on one continuation makes the states downstream of
that continuation more likely and the states downstream of its abandoned siblings
less likely; the diversity those siblings carried leaves the trajectory distribution
even though no conditional distribution anywhere became more peaked.

Our first contribution is to show that the local channel is not merely blind to the
second term but \emph{anti-correlated} with it. The local estimator factors into the
advantage, a saturation factor, and the deviation of the sampled token's surprisal
from the mean surprisal --- and that last factor has mean zero by the definition of
entropy. At a branch point where several continuations are close to equiprobable it
vanishes identically. The vanishing is not an estimator failure: a uniform
distribution maximizes entropy, so the first-order change is genuinely zero and the
damage is second order (Section~\ref{sec:blind}). Meanwhile the visitation term's
prefactor is $1-\pit(y_t \mid s_t)$, which is maximal precisely at those flat
states. The local channel therefore spends its budget where the conditional entropy
is fragile --- rare tokens sampled from confident distributions --- and spends
nothing where the trajectory entropy is most at risk.

Our second contribution is an estimator for the missing channel that requires no
additional rollouts. Eliminating the non-causal cross terms turns the visitation
term into a REINFORCE estimator whose reward-to-go is the sum of the
\emph{remaining} conditional entropies, $\Htogo$; subtracting any function of the
state leaves it unbiased. We estimate $\Htogo$ with parallel
multi-token-prediction (MTP) heads read off the hidden state the policy already
computes, and we obtain the state-only baseline from rollouts that share a prefix.
Plain group-relative rollouts almost never share prefixes, so we schedule the
sharing: a tree-structured sampler branches a single root at fixed depths, and the
resulting sibling sets give the baseline an actual support
(Section~\ref{sec:tree}).

The resulting per-token score $\AH$ is a deviation from a sibling mean and therefore
sums to zero over each sibling set. This has a consequence we state up front because
it determines what evidence can count: the method \emph{cannot} raise the average
weight at branch points, only spread the weights of siblings apart. ``Total entropy
went up'' is therefore not a prediction of this method, and an experiment looking
for it is testing a different hypothesis (Section~\ref{sec:zerosum}).

We make the following contributions.

\begin{itemize}
\item An exact two-channel decomposition of the trajectory-entropy gradient, and an
      argument that the channel captured by current token-level entropy-control
      methods is anti-correlated with the channel they miss
      (Sections~\ref{sec:channels}--\ref{sec:blind}).
\item A derivation of the missing channel into a per-token advantage $\AH$ with no
      approximation up to the point of estimating $\Htogo$, together with an explicit
      ledger of the approximations that estimation requires, the sign of each one's
      error, and the mechanism that absorbs it
      (Sections~\ref{sec:ahdef}--\ref{sec:ledger}).
\item A tree-structured rollout that supplies the sibling-prefix baseline, and a
      bound showing what fraction of positions can carry a non-zero $\AH$ under
      \emph{any} sampler (Section~\ref{sec:tree}).
\item A five-arm experimental design that separates the contribution of the
      forecast's value from the contribution of the apparatus that delivers it, by
      running the full pipeline twice with the forecast's value destroyed in two
      different ways (Section~\ref{sec:setup}).
\end{itemize}

\section{Related Work}
\label{sec:related}

\paragraph{Entropy collapse in RLVR.}
Entropy decay under policy-gradient fine-tuning is well documented, and a range of
interventions add an explicit entropy term to the objective or constrain the update
so that high-entropy positions are preserved. Recent analyses report that naive
entropy regularization does not prevent the collapse and that continuing to train
under it degrades accuracy \citep{rethinkingentropyreg}, which is consistent with
our reading: the regularizer acts on the local channel, and the local channel is
not where trajectory diversity is being lost.

\paragraph{Token-level entropy control.}
The method we build on re-weights the clipped surrogate loss by a per-token
first-order estimate of the entropy change that token's update induces
\citep{steer}. That work observes that earlier interventions are ``indirect'' and
that their effect is therefore limited and can fail. Our decomposition supplies a
mechanism for that observation: the estimate is a first-order estimate of the local
channel only, and it is exactly zero at the states where the visitation channel is
largest.

\paragraph{Which tokens carry the gradient.}
\citet{eightytwenty} show that restricting policy-gradient updates to a small
minority of high-entropy ``forking'' tokens matches full-gradient training. That
result establishes that the identity of the tokens, not the aggregate entropy,
is what matters, and it motivates a per-token treatment; it does not, however,
distinguish the two channels, since forking tokens are identified by their local
entropy. Our analysis predicts that a local criterion will select a systematically
different set from the one the visitation channel would select.

\paragraph{Local stochasticity without global diversity.}
\citet{invisibleleash} report that token-level entropy can remain high while the
diversity of final answers falls. This is precisely the signature our decomposition
predicts when the visitation channel is unmanaged: per-state entropies are
preserved --- the local channel is doing its job --- while the trajectory
distribution concentrates.

\paragraph{Multi-token prediction.}
We use Medusa-style parallel prediction heads \citep{medusa} not for their intended
purpose (speculative decoding) but as an entropy forecaster: we need the entropy of
the distribution over a token $k$ steps ahead, not a good sample of that token.
Section~\ref{sec:approx} shows that this choice of estimator has a signed,
monotone error which we characterize and correct.

\section{Preliminaries}
\label{sec:prelim}

\paragraph{Notation.}
A prompt is $x$, a response is $y = y_{1:T}$, and the state before emitting token
$t$ is $s_t = (x, y_{<t})$. We write $H(\pit(\cdot \mid s_t))$ for the entropy of
the conditional distribution at a single position (lowercase, per-token) and
$\Htraj(\pit \mid x)$ for the entropy of the distribution over whole responses
(uppercase, per-trajectory). The surprisal of a token $a$ is
$I_a := -\log \pit(a \mid s_t)$, so that $H = \E_{a \sim \pit}[I_a]$.

\paragraph{Group-relative policy optimization.}
We optimize with GRPO. For each prompt, $n$ responses are sampled and scored, and
each response receives the group-standardized advantage
$A_i = (r_i - \mathrm{mean}_j r_j) / \mathrm{std}_j r_j$, broadcast to every token
of that response. Replacing a learned value function with the group mean is what
makes GRPO cheaper than PPO; it also means that a group whose members all receive
the same reward contributes no gradient.

\paragraph{Token-level entropy control.}
The baseline intervention leaves the advantage untouched and instead multiplies each
token's contribution to the clipped surrogate by a weight
$\alpha_t \in [w_{\min}, w_{\max}]$ derived from a first-order estimate $\Om_t$ of
the entropy change that token's update induces. With $w_{\max} = 1$, the weight can
only attenuate. Writing $A_u$ for the advantage, $\pi_u := \pit(y_u \mid s_u)$ for
the probability of the sampled token, and $H_u$ for the local entropy, the estimator
used in that line of work is, up to a positive constant,
%
\begin{equation}
\Om_u \;=\; A_u\,\pi_u (1-\pi_u)\,\bigl(\log \pi_u + H_u\bigr).
\label{eq:omega}
\end{equation}
%
The policy loss becomes $\sum_u \alpha_u \cdot \mathrm{PPO\text{-}clip}(A_u, r_u)$,
where the weights are a monotone map of $\Om$ over the batch.

\section{Two Channels}
\label{sec:channels}

\subsection{The decomposition is exact}

Applying the entropy chain rule to $\pit(y \mid x) = \prod_t \pit(y_t \mid s_t)$
gives an identity, not an approximation:
%
\begin{equation}
\Htraj(\pit \mid x)
= \E_{y \sim \pit}\Bigl[\sum_{t=1}^{T} H\bigl(\pit(\cdot \mid s_t)\bigr)\Bigr].
\label{eq:chain}
\end{equation}
%
The left-hand side is what we would like to protect; the right-hand side is what
instrumentation can measure. The entire argument of this paper follows from the
observation that the measure $\E_{y \sim \pit}$ depends on $\theta$ as well.

Set $f_\theta(y) := \sum_t H(\pit(\cdot \mid s_t(y)))$, so that
$\Htraj = \E_{y \sim \pit}[f_\theta(y)]$. Both the integrand and the measure depend
on $\theta$, so the product rule --- with the score-function identity restoring the
second term to an expectation --- yields
%
\begin{equation}
\begin{split}
\gradt \Htraj
&= \underbrace{\E\bigl[\gradt f_\theta(y)\bigr]}_{(\mathrm{L})\ \text{local}}\\
&\quad + \underbrace{\E\bigl[f_\theta(y)\,\gradt \log \pit(y \mid x)\bigr]}_{(\mathrm{V})\ \text{visitation}}.
\end{split}
\label{eq:twochannel}
\end{equation}
%
\textbf{(L), the local channel}, holds the visited states fixed and asks how the
conditional distribution at each one changes. Equation~\eqref{eq:omega} is a
first-order estimate of exactly this.

\textbf{(V), the visitation channel}, holds the per-state entropies fixed and asks
how the probability of reaching each state changes. An update that concentrates mass
on one continuation makes the sub-tree below it more likely and the sub-trees below
its siblings less likely. Nothing in Eq.~\eqref{eq:omega} can see this: $\Om$ reads
the conditional distribution at a position and never reads the probability of
arriving there.

The two terms can have opposite signs. A method that drives (L) to zero has not
thereby bounded $\gradt \Htraj$.

\subsection{Why the local channel cannot be repaired}
\label{sec:blind}

A natural objection is that the two channels might coincide in practice: branch
points are positions where the policy is uncertain, so they are high-entropy
positions, so $\Om$ should be large there and the local term should protect them
automatically. The premise is correct --- in warm-up rollouts on our model, branch
diversity correlates with local entropy at Spearman $\rho = 0.65$, and with the
entropy-to-go at $\rho = 0.82$. The inference from it fails, in three steps.

\paragraph{Step 1: $\Om$ reads a deviation, not a level.}
Rewriting Eq.~\eqref{eq:omega} in terms of surprisal gives
%
\begin{equation}
\Om \;=\; A\,\pi_a(1-\pi_a)\,\bigl(H - I_a\bigr),
\label{eq:omega-surprisal}
\end{equation}
%
and $H = \E_{a\sim\pi}[I_a]$ by definition. The factor $H - I_a$ is therefore a
mean-zero fluctuation: what sets the magnitude of $|\Om|$ is the \emph{spread} of
surprisal, not the \emph{height} of the entropy.

\paragraph{Step 2: at a uniform branch point $\Om \equiv 0$.}
If the distribution is uniform over $k$ candidates then $I_a = \log k = H$ for every
$a$, so $\Om$ vanishes identically. The positions that are most branch-like --- several
continuations at comparable probability --- are exactly where the local term
contributes nothing.

\paragraph{Step 3: and that zero is the correct answer.}
This is the step that closes the argument. A uniform distribution maximizes entropy,
so $\partial H / \partial z_a = 0$ for every logit $z_a$; $\Om$ is not wrong, it is
right. Perturbing one logit of a uniform distribution by $\delta$ and writing
$S := e^{\delta} + k - 1$ gives the closed form
%
\begin{equation}
H(\delta) = \log S - \frac{\delta e^{\delta}}{S},
\qquad
H''(0) = -\,\frac{k-1}{k^{2}},
\label{eq:secondorder}
\end{equation}
%
with $H'(\delta)$ of sign $-\operatorname{sign}(\delta)$: any push in any direction
lowers the entropy, but the loss is entirely second order. For a $10$-way uniform
distribution and $\delta = 0.1$ the actual entropy loss is $4.7\times10^{-4}$,
roughly two orders of magnitude below that of a peaked distribution under the same
perturbation. \emph{As far as local entropy is concerned, uniform branch points are
genuinely safe}, and a local estimator that allocates them no protection is
allocating correctly for its own objective. No refinement of the local term
recovers the missing signal, because the missing signal is not local.

\paragraph{The two channels are anti-correlated.}
The visitation term's prefactor, derived in Section~\ref{sec:weights}, is
$1 - \pi_a$, which approaches $1$ as the distribution flattens --- and the
probability that a group of $n$ samples actually diverges at a position approaches
$1$ in the same limit. Table~\ref{tab:anticorr} evaluates both channels on a family
of distributions that interpolates between confident and flat. $|\Om|$ is dominated
by positions where the model is fairly confident but a tail token was sampled; at
those positions siblings diverge only about half the time. Where divergence is
certain, the local term is identically zero and the visitation prefactor is maximal.

\begin{table}[t]
\centering
\footnotesize
\setlength{\tabcolsep}{3.5pt}
\begin{tabular}{lcccc}
\toprule
distribution & $H$ & $P(\mathrm{split})$ & $\E|\Om|/|A|$ & $\E[1{-}\pi_a]$ \\
\midrule
$q{=}.99$          & 0.10 & 0.08 & 0.092 & 0.020 \\
$q{=}.9$, tail     & 1.25 & 0.57 & \textbf{1.129} & 0.190 \\
$q{=}.6$, $m{=}3$  & 1.11 & 0.98 & 0.457 & 0.587 \\
unif.\ $k{=}10$    & 2.30 & \textbf{1.00} & \textbf{0.000} & \textbf{0.900} \\
unif.\ $k{=}100$   & 4.61 & \textbf{1.00} & \textbf{0.000} & \textbf{0.990} \\
\bottomrule
\end{tabular}
\caption{The local channel and the visitation channel are ordered oppositely.
Distribution family: top token has mass $q$, the remaining $1-q$ is spread uniformly
over $m$ tokens. $P(\text{split})$ is the probability that $n{=}8$ samples do not all
agree. $\E|\Om|$ is taken over the token actually sampled at that position.}
\label{tab:anticorr}
\end{table}

The statement we can therefore make is stronger than ``the local term is
incomplete'': \emph{the local term is anti-correlated with the term it omits}, and
the omission cannot be fixed within the local channel.

\section{Estimating the Visitation Channel}
\label{sec:ahdef}

\subsection{Causal elimination gives a reward-to-go}

Substituting $\gradt \log \pit(y \mid x) = \sum_u \gradt \log \pit(y_u \mid s_u)$
into (V) produces a double sum over $(t, u)$. For $u \ge t$ the entropy
$H(\pit(\cdot \mid s_t))$ is a function of $y_{<t}$, hence measurable with respect
to $y_{<u}$, and the inner conditional expectation vanishes:
%
\begin{equation*}
\begin{split}
\E\bigl[\gradt \log \pit(y_u \mid s_u) \mid y_{<u}\bigr]&\\
= \gradt \!\sum_a \pit(a \mid s_u) &= 0 .
\end{split}
\end{equation*}
%
Only $u < t$ survives, and the surviving inner sum is the entropy still to come:
%
\begin{equation}
(\mathrm{V}) = \E\Bigl[\sum_u \gradt \log \pit(y_u \mid s_u)\,
\Htogo(s_u \oplus y_u)\Bigr],
\label{eq:togo}
\end{equation}
%
with $\Htogo(s_u \oplus y_u) := \sum_{t > u} H(\pit(\cdot \mid s_t))$.

We stress that $\Htogo$ was not posited. Equation~\eqref{eq:togo} is a REINFORCE
estimator whose reward is ``the conditional entropy remaining after this token'',
and $\Htogo$ is that reward-to-go.

\subsection{The sibling-prefix baseline}

Any function $b(s_u)$ of the state alone leaves Eq.~\eqref{eq:togo} unbiased, by the
same argument. Subtracting one gives the per-token advantage
%
\begin{equation}
\AH(s_u, y_u) := \Htogo(s_u \oplus y_u) - \bar{H}_{\mathrm{togo}}(s_u),
\label{eq:ah}
\end{equation}
%
and up to this point \emph{no approximation has been made}.

The baseline we use is the mean of $\Htogo$ over the rollouts in the same prompt
group whose prefixes $y_{<u}$ agree exactly. Those rollouts share the state $s_u$ by
definition, so their mean is a function of $s_u$ alone and the baseline is
admissible. The condition is agreement on $y_{<u}$ and not on $y_{\le u}$: the
rollouts must diverge \emph{at} $u$ for the difference to be a branch score.

Two properties follow immediately and matter later. First, the implementation
includes a rollout in its own sibling set, so $b$ depends weakly on $y_u$ and the
estimator is not strictly leave-one-out; the resulting bias is $O(1/m)$ in the
number $m$ of surviving siblings, the same property GRPO's group-mean advantage
has. Second, when $m = 1$ the baseline equals the rollout's own value and
$\AH \equiv 0$, so a position with no siblings is silent rather than noisy.

\section{Four Approximations}
\label{sec:approx}

Equation~\eqref{eq:ah} is exact but not computable: $\Htogo$ sums conditional
entropies at states the update has not visited. Three approximations make it
computable and a fourth bounds its influence.

\subsection{A1: horizon truncation and discounting}

We truncate the sum to $\kappa$ steps and discount:
%
\begin{equation}
\Htogo(s_u \oplus y_u) \approx
\sum_{k=1}^{\kappa} \gamma_H^{\,k}\, H\bigl(\pi(\cdot \mid s_{u+k})\bigr).
\label{eq:a1}
\end{equation}
%
The discount serves two purposes: distant terms carry more forecast error (A2) and
therefore deserve less trust, and the sum stays bounded independently of $\kappa$.

What this discards --- the tail $\sum_{k > \kappa}$ and a fraction of the retained
terms --- is largely harmless for a reason specific to our use. Equation~\eqref{eq:ah}
consumes only the \emph{difference} between siblings. Any bias common to rollouts
sharing $s_u$ cancels in that difference, and truncation bias is common to first
order.

\subsection{A2: forecasting the unrealized future}

Even Eq.~\eqref{eq:a1} requires the future states $s_{u+k}$. At update time they do
not exist: rollouts came from $\pi_{\mathrm{old}}$, each state has exactly one
sampled continuation, and re-rolling is prohibitive. We replace the realized future
with a parallel forecast from the current hidden state $h_u$:
%
\begin{equation}
H\bigl(\pi(\cdot \mid s_{u+k})\bigr) \;\longrightarrow\;
H\bigl(p^{(k)}_{\mathrm{MTP}}(\cdot \mid h_u)\bigr).
\label{eq:a2}
\end{equation}
%
Head $k$ is a two-layer residual MLP applied to $h_u$, followed by the policy's own
unembedding; the heads own no vocabulary parameters. Their last layer is
zero-initialized, so before any training every head reproduces the policy's own
next-token distribution and the forecast starts from a sane value. In the RL path the
logits are materialized in chunks and reduced to an entropy immediately, so the
$[\kappa, B, T, V]$ tensor never exists.

\paragraph{The error of this substitution has a known sign.}
Conditioning on the intervening tokens is exactly what the forecast omits. Fixing
$S_u$, the $k$-th term of Eq.~\eqref{eq:a1} has expectation
$H(Y_{u+k} \mid Y_{u+1:u+k-1}, S_u)$, whereas the MTP head models the
\emph{marginal} $p(Y_{u+k} \mid S_u)$. Since conditioning reduces entropy,
%
\begin{equation}
\begin{split}
H(Y_{u+k} \mid S_u) &- \E\bigl[H(\pi(\cdot \mid S_{u+k-1}))\bigr]\\
&= I\bigl(Y_{u+k};\, Y_{u+1:u+k-1} \mid S_u\bigr)\\
&\ge 0 .
\end{split}
\label{eq:mi}
\end{equation}
%
The forecast \emph{always over-estimates}, and the excess is precisely the mutual
information with the intervening tokens, which grows monotonically with $k$.

Equation~\eqref{eq:mi} is not merely a bound: it is the reason the marginal is the
quantity we want. It is the fan-out of the continuation --- how many ways the next
$k$ tokens can unfold --- rather than the residual uncertainty along the one path
that happened to be sampled.

\subsection{A3: per-head affine calibration}

Equation~\eqref{eq:mi} predicts a per-head, monotone inflation, so we absorb it with
a per-head affine correction fitted by least squares against oracle entropies on
warm-up rollouts:
%
\begin{equation}
\hat{H}_k = a_k\, H\bigl(p^{(k)}_{\mathrm{MTP}}(\cdot \mid h_u)\bigr) + b_k .
\label{eq:a3}
\end{equation}
%
The fitted coefficients confirm the prediction. On our model the first head has
$a_1 = 1.03$ --- a one-step forecast has no intervening tokens, so the mutual
information in Eq.~\eqref{eq:mi} is zero and the approximation is exact --- while
$a_k$ drops to $0.10$--$0.19$ for $k \ge 2$. Calibration switches the unreliable
heads off on its own.

Combining A1--A3 gives the forecast that is actually computed:
%
\begin{equation}
\begin{split}
&\hat{H}_{\mathrm{togo}}(s_u \oplus y_u) =\\
&\qquad \max\Bigl(0,\ \sum_{k=1}^{\kappa} \gamma_H^{\,k}\bigl(a_k H_k(h_u) + b_k\bigr)\Bigr),
\end{split}
\label{eq:forecast}
\end{equation}
%
the clamp encoding the fact that an entropy cannot be negative. With
$\kappa = 2$ and $\gamma_H = 0.7$ this instantiates to
$\hat{H}_{\mathrm{togo}} = 0.72\,H_1 + 0.06\,H_2 + 0.11$: the second head
contributes about $8\%$, so the effective horizon is close to one step. We regard
this as a consequence of Eq.~\eqref{eq:mi} rather than a tuning artifact, and note
it as a limitation --- a forecaster whose error did not grow with $k$ would leave a
longer horizon usable.

\paragraph{Alignment.}
One detail decides whether the signal is a branch score at all. The hidden state at
position $u$ is conditioned on $s_u$ and predicts $y_u$; it does not know $y_u$. The
quantity we need, $\Htogo(s_u \oplus y_u)$, must therefore be read from position
$u+1$ and shifted back by one. Omitting the shift measures the value of the state
\emph{before} the branch rather than the value of the branch taken.

\section{From Gradient to Token Weights}
\label{sec:weights}

Equation~\eqref{eq:ah} is a term to add to a gradient, but methods in this family do
not add gradients --- they re-weight a loss. What is needed is a first-order
prediction of how much one update, acting through token $u$, changes the trajectory
entropy.

\paragraph{How far one step moves a logit.}
With $w_u = \mathrm{clip}(r_u, 1-\epsilon_{\mathrm{lo}}, 1+\epsilon_{\mathrm{hi}})A_u$
the coefficient the policy loss actually applies, a single step moves only the
sampled token's logit, $\Delta z_{u,a} = \eta\, w_u\, \delta_{a, y_u} + O(\eta^2)$.
The softmax Jacobian $\partial \log \pi_a / \partial z_a = 1 - \pi_a$ then gives
%
\begin{equation}
\widehat{\Delta \log \pi}_u = \eta\, w_u\,(1 - \pi_u),
\label{eq:dlogpi}
\end{equation}
%
which vanishes as $\pi_u \to 1$: a token that already holds the mass cannot be given
more. This is the prefactor whose anti-correlation with $|\Om|$ Table~\ref{tab:anticorr}
records.

\paragraph{The visitation-channel prediction.}
Because $\gradt \log \pi_u \cdot \Delta\theta = \Delta \log \pi_u$, substituting into
Eq.~\eqref{eq:ah} gives the per-token quantity
%
\begin{equation}
\begin{split}
\mathrm{visit}_u := \widehat{\Delta \log \pi}_u \cdot
\mathrm{clip}\bigl(\AH(s_u,y_u), -c, +c\bigr).
\end{split}
\label{eq:visit}
\end{equation}
%
\textbf{A4} is the clip, and it is applied to $\AH$ \emph{before} the product rather
than after, so that one token's influence is bounded by $|\eta w_u| \cdot c$.
Clipping the product instead would let forecast error be amplified without bound at
tokens with large $w_u$.

\paragraph{Sign quadrants.}
Equation~\eqref{eq:visit} defines the collapse channel precisely.
When $\widehat{\Delta \log \pi}_u > 0$ and $\AH > 0$, mass moves toward a branch
richer than its siblings and trajectory entropy rises. When
$\widehat{\Delta \log \pi}_u > 0$ and $\AH < 0$, mass moves into a dead end and the
diversity carried by the abandoned siblings is destroyed --- this is the collapse
path that (L) is structurally unable to see. Negative
$\widehat{\Delta \log \pi}_u$ reverses both readings.

\section{Combining the Channels}
\label{sec:combine}

The two channels enter the same weight:
%
\begin{equation}
\Omt_u = \mathrm{norm}(\Om_u) + \lambda \cdot \mathrm{norm}(\mathrm{visit}_u),
\label{eq:combine}
\end{equation}
%
after which $\Omt$ is passed to the unmodified min--max mapping of the base method,
producing $\alpha_u \in [w_{\min}, w_{\max}]$. We emphasize that the hook is
inserted where $\Om$ is finalized and that nothing downstream is changed; when
$\lambda = 0$ the pipeline is bitwise identical to the base method, which we
enforce with a regression test.

Three choices in Eq.~\eqref{eq:combine} deserve comment. Normalization is required
even though both terms carry units of nats, because $\Om$ contains a factor
$1/\pi_{\mathrm{old}}$ and is heavy-tailed. Because normalization removes scale, the
step size $\eta$ in Eq.~\eqref{eq:dlogpi} only ever appears multiplied by $\lambda$;
we fix $\eta = 1$ and sweep $\lambda$ alone. Finally, $\mathrm{norm}$ is an RMS
rescaling rather than a $z$-score, since subtracting a mean would break the
$\lambda = 0$ equivalence.

\section{Tree-Structured Rollouts}
\label{sec:tree}

\subsection{Why plain rollouts cannot support $\AH$}

$\hat{H}_{\mathrm{togo}}$ is a deterministic function of the causal prefix, so two
rollouts sharing $y_{\le u}$ produce identical forecasts. Combined with
Eq.~\eqref{eq:ah} this gives a sharp characterization:
%
\begin{equation}
\begin{split}
\AH[i, u] \neq 0 \iff\ &\text{siblings share } y_{<u}\\
&\text{but diverge at } y_u,
\end{split}
\label{eq:branchonly}
\end{equation}
%
that is, $u$ is a branch point of the group. Under independent sampling from a
sharpened policy, group members diverge at their first token and never share a
prefix again, so $m = 1$ almost everywhere and $\AH \equiv 0$: the visitation
channel is not merely weak, it is identically absent.

We therefore schedule the sharing. A tree sampler expands a single root and cuts it
at fixed depths $d_1 < d_2 < d_3$ with branching factor $2$, yielding $n = 8$ leaves
whose sibling structure is known in advance. The cuts create the only positions at
which Eq.~\eqref{eq:branchonly} can fire.

\subsection{A ceiling on the support}

It is tempting to treat the fraction of positions with $\AH \neq 0$ as a quantity to
be maximized. It is not, and the reason is combinatorial rather than empirical. Two
quantities behave very differently under the tree:

\begin{itemize}
\item the fraction of positions where \emph{some} other rollout shares $y_{<u}$ is
      moved by the cuts, and equals $d_L/T$ by construction;
\item the fraction where $\AH \neq 0$ is \emph{not}, because by
      Eq.~\eqref{eq:branchonly} it counts branch points, and $n$ sequences form a
      trie with at most $n-1$ internal branching nodes.
\end{itemize}

Since the one-position alignment shift widens each branch node into two columns, the
fraction of positions carrying a non-zero $\AH$ is bounded by $2(n-1)/T$ for
\emph{any} sampler whatsoever. With $n = 8$ and responses of $\sim 10^3$ tokens this
ceiling is on the order of a percent. The tree's role is to make the branch points
exist and to place them at controlled depths, not to make them numerous.

\subsection{The correction cannot move the mean}
\label{sec:zerosum}

$\AH$ is a deviation from a sibling mean, so on each sibling set
%
\begin{equation}
\sum_{j \in \mathrm{sib}(s_u)} \AH\bigl(s_u, y^{(j)}_u\bigr) = 0 .
\label{eq:zerosum}
\end{equation}
%
The correction in Eq.~\eqref{eq:combine} therefore cannot change the average weight
assigned to branch tokens. It can only spread siblings apart. This is not a defect
to be engineered away: an unbiased estimator requires a baseline, and a baseline
forces the mean to zero.

The consequence for evaluation is direct, and we state it before presenting any
experiment. ``Mean entropy at branch points increased'' is \emph{not} a prediction of
this method and cannot serve as evidence that it works; the corresponding prediction
is that the \emph{variance across siblings} widens. A method designed to raise the
mean is a different and weaker hypothesis --- it does not need a forecast at all,
only a support --- and we include it below as a control arm.

\section{The Approximation Ledger}
\label{sec:ledger}

Table~\ref{tab:ledger} collects every departure from Eq.~\eqref{eq:twochannel}.

\begin{table*}[t]
\centering
\footnotesize
\setlength{\tabcolsep}{4pt}
\begin{tabular}{@{}lp{2.6cm}p{4.0cm}p{2.3cm}p{4.1cm}@{}}
\toprule
& approximation & what it discards & sign / bound & mitigation \\
\midrule
--- & Eqs.~\eqref{eq:chain}--\eqref{eq:ah} & nothing --- identities & --- & --- \\
A1 & horizon $\kappa$, discount $\gamma_H$ & the tail $\sum_{k>\kappa}$, and a fraction of the retained terms & under-estimate & the common component cancels in the sibling difference \\
A2 & realized future $\to$ MTP marginal & mutual information with intervening tokens, Eq.~\eqref{eq:mi} & \textbf{over-estimate, monotone in $k$} & A3 \\
A3 & per-head affine calibration & non-affine residual of A2 & least-squares residual & oracle-forecast control arm \\
A4 & clip $\AH$ to $\pm c$ & magnitude of extreme branch scores & bounded by $c$ & saturation rate is logged \\
B1 & first order in $\eta$ & $O(\eta^2)$ & --- & small learning rate \\
B2 & only the sampled logit moves & cross-position effects of weight sharing & --- & \emph{inherited from the base method} \\
B3 & rollouts come from $\pi_{\mathrm{old}}$ & mismatch with the $\pit$ expectation & --- & importance ratio restores $r$ \\
B4 & self included in sibling set & $O(1/m)$ relative to leave-one-out & --- & same property as the GRPO advantage \\
\bottomrule
\end{tabular}
\caption{Every approximation between the exact decomposition and the implemented
weight. A1--A4 are introduced by this work; B1--B4 are shared with the token-level
entropy-control method we extend.}
\label{tab:ledger}
\end{table*}

The strongest assumptions are A2 and B2. A2 has a known sign and a known monotonicity
(Eq.~\ref{eq:mi}), which is what makes the affine correction in A3 an appropriate
remedy rather than a curve fit. B2 --- that updating one token does not change the
conditional distribution at other positions --- is false for a network with shared
parameters and has no corrective mechanism. It is not a risk this work introduces,
since the base method already relies on it, but adding a second channel adds a second
copy of its error.

\section{Experimental Setup}
\label{sec:setup}

\subsection{Model, data, and optimization}

All runs fine-tune the same base model on the same data with the same optimizer
settings, and differ only in the intervention. We use a $1.5$B-parameter
mathematics-specialized model, train on a $17$k-problem competition-mathematics
set, and validate on AIME 2024 ($30$ problems) with $32$ samples per problem.
With a training batch of $512$ prompts, $110$ optimization steps correspond to
about $3.2$ epochs.

Optimization uses GRPO with $n = 8$ responses per prompt, a constant learning rate
of $1\mathrm{e}{-6}$ with no warm-up (so that the schedule is independent of the
declared total step count), a maximum response length of $3072$ tokens, no entropy
bonus, and no KL loss. Rollouts are sampled at temperature $1.0$ with no nucleus
truncation; validation uses temperature $1.0$ and $\mathrm{top\text{-}}p = 0.7$.
Every arm uses the same data seed and the same initial checkpoint, which we verify by
requiring the step-$0$ validation scores to agree.

\subsection{Method hyperparameters}

The forecaster uses $\kappa = 2$ heads with discount $\gamma_H = 0.7$ and the
per-head calibration of Eq.~\eqref{eq:a3}. The combination uses $\lambda = 0.25$,
RMS normalization, clip $c = 1.0$, the sibling-prefix baseline, and the base
method's min--max mapping into $[w_{\min}, w_{\max}] = [0.7, 1.0]$ --- so the weight
can only attenuate, never amplify. The tree sampler uses a single root, cuts at
depths $\{64, 192, 384\}$, and branching factor $2$ at each cut, giving $n = 8$
leaves.

\subsection{Five arms}

The design isolates one question: does the \emph{value} of the forecast matter, or
only the apparatus that carries it? A comparison between the base method and the
full method changes two things at once ($\lambda$ and the sampler), so it cannot
answer this. We therefore add two arms that keep the entire apparatus --- the tree,
the forecaster, the support, the weight mapping --- and destroy only the information
in $\AH$.

\begin{table}[t]
\centering
\footnotesize
\setlength{\tabcolsep}{4pt}
\begin{tabular}{@{}llll@{}}
\toprule
arm & $\lambda$ & rollout & signal in the weight \\
\midrule
GRPO      & ---   & plain & none \\
STEER     & $0$   & plain & local $\Om$ only \\
uniform   & $.25$ & tree  & support only \\
permuted  & $.25$ & tree  & $\AH$ shuffled \\
\textbf{STEER-F} & $.25$ & tree & $\AH$ as derived \\
\bottomrule
\end{tabular}
\caption{The five arms. \textsc{uniform} and \textsc{permuted} run the full pipeline
with the forecast's value destroyed in two different ways, so any gain they show is
attributable to the apparatus rather than to the visitation signal.}
\label{tab:arms}
\end{table}

\textsc{uniform} applies the same attenuation to every position in the support,
discarding $\AH$'s value but keeping its support --- this is the implementation of
the ``raise the mean'' hypothesis that Section~\ref{sec:zerosum} distinguishes from
ours. \textsc{permuted} keeps the multiset of $\AH$ values within each sibling set and
shuffles the assignment, so the support, the magnitudes, and Eq.~\eqref{eq:zerosum}
all survive while the pairing between a branch and its score is destroyed. Together
they bracket the claim: if the visitation signal is doing the work, both controls
should behave like the no-forecast arms.

\subsection{Metrics}

We report $\mathrm{acc@}32$ (mean accuracy over $32$ samples) and $\mathrm{maj@}32$
(accuracy of the majority vote over the same $32$ samples), and their difference
%
\begin{equation*}
\mathrm{uplift} := \mathrm{maj@}32 - \mathrm{acc@}32,
\end{equation*}
%
which measures how much better a set of $32$ samples is than its average member ---
that is, whether correct answers concentrate while incorrect ones scatter. We use
uplift rather than $\mathrm{pass@}32$ because the latter does not separate the arms.

Because single validation points are noisy on a $30$-problem set, the primary
statistic is the mean over the converged window, steps $40$--$110$ ($8$ validation
points per arm). We quote differences in units of the across-problem standard error,
$\mathrm{SE} = \mathrm{std@}32 / \sqrt{30} \approx 0.027$, which is the correct scale
for generalization to new problems. We separately note that the pure re-sampling
noise of the protocol --- estimated from two validations of identical weights --- is
about $0.004$ in accuracy, an order of magnitude smaller; the two scales answer
different questions and we are explicit about which is in use.

One instrumentation caveat propagates into every entropy number we report. The
logged \texttt{actor/entropy} is aggregated with a token-mean, so it is
$\Htraj$ divided by the mean response length, not $\Htraj$. Since response lengths
differ by under $5\%$ across arms, we report the raw quantity and, where the
distinction could matter, the length-corrected product as well.

\section{Results}
\label{sec:results}

\TODO{This section is deliberately left as scaffolding. The training logs for four
of the five arms are available; the GRPO arm's log and the multi-benchmark
evaluation are pending, and the numbers will be filled in once both are verified
from source.}

\subsection{Main comparison}
\label{sec:res:main}

\TODO{Table: per-arm plateau means over steps 40--110 --- acc@32, maj@32, uplift,
entropy, mean response length, length-corrected entropy. Five rows, one per arm of
Table~\ref{tab:arms}.}

\subsection{Does the forecast's value matter?}
\label{sec:res:contrasts}

\TODO{Table of pairwise contrasts in SE units: STEER-F versus each of GRPO, STEER,
uniform, permuted, and the three null contrasts among the control arms. The
prediction of Section~\ref{sec:setup} is that the first group separates and the
second does not.}

\subsection{Decomposing the gain}
\label{sec:res:decomp}

\TODO{Split the maj@32 gain into the part attributable to per-sample accuracy and
the part attributable to uplift, and split the total gain into the apparatus
component (permuted minus base) and the information component (STEER-F minus
permuted).}

\subsection{Entropy is not the channel}
\label{sec:res:entropy}

\TODO{Table contrasting entropy gaps against accuracy gaps across arms, testing
whether arms with matched aggregate entropy differ in accuracy. Section~\ref{sec:zerosum}
predicts that they can, and that aggregate entropy therefore does not order the
arms.}

\subsection{Mechanism diagnostics}
\label{sec:res:diag}

\TODO{Support fraction, mean sibling-set size, and branch-point fraction per step,
against the $2(n-1)/T$ ceiling of Section~\ref{sec:tree}; the two-term balance
$\lambda\|\mathrm{visit}\| / \|\Om\|$; and the $\AH$ clip saturation rate.}

\subsection{Cost}
\label{sec:res:cost}

\TODO{Seconds per step per arm and the resulting wall-clock overhead, attributed to
the MTP forward pass.}

\section{Conclusion}

The gradient of trajectory entropy has two terms, and the token-level entropy
controls now in use estimate one of them. We showed that the omission is structural:
the estimated term vanishes exactly where the omitted term is largest, and it
vanishes for a correct reason, so the gap cannot be closed by improving the
estimator. We then showed that the omitted term is a REINFORCE estimator whose
reward-to-go is the remaining conditional entropy, that a parallel forecaster
estimates it for the price of one forward pass with an error whose sign and
monotonicity are known in closed form, and that a tree-structured rollout supplies
the baseline it requires. The resulting correction is mean-preserving by
construction: it does not slow entropy collapse, it redirects it.

\section*{Limitations}

\paragraph{A single validation benchmark.}
Every claim rests on AIME 2024, which has $30$ problems. At $\mathrm{SE} \approx
0.027$, differences of interest are a fraction of one problem, and a multi-benchmark
evaluation is required before the direction of any gap can be called robust. We
report the standard error explicitly rather than significance stars for this reason.

\paragraph{A single seed.}
All arms use one data seed. The standard error we quote is an across-problem
quantity and does not bound run-to-run variation, which for these runs is
uncontrolled because rollout sampling is not seeded across processes. A re-run of
the same configuration is a second draw, not a restoration.

\paragraph{The effective horizon is one step.}
The calibration in Eq.~\eqref{eq:a3} suppresses heads beyond the first, so
$\hat{H}_{\mathrm{togo}}$ is dominated by a single step ahead. This is the predicted
behavior of Eq.~\eqref{eq:mi} rather than a bug, but it means the ``future'' in the
visitation term is a short future, and a forecaster whose error did not grow with
$k$ would test the idea more sharply. A sequential forecaster is the obvious
comparison and we have not run it.

\paragraph{The support is small by construction.}
Section~\ref{sec:tree} bounds the fraction of positions carrying a non-zero $\AH$ at
$2(n-1)/T$, on the order of a percent for our settings. Whether the effect scales
with $n$, or with a branching schedule chosen adaptively rather than at fixed
depths, is untested.

\paragraph{An assumption with no remedy.}
Approximation B2 --- that a token's update does not change the conditional
distribution at other positions --- is false in a network with shared parameters.
It is inherited from the base method rather than introduced here, but our
construction adds a second estimator that relies on it.

\paragraph{One prediction of the analysis is unmeasured.}
Section~\ref{sec:blind} argues that the two channels are anti-correlated. The
prefactor half of that claim is analytic and the $\AH$ half is inferred from a
correlation with branch diversity; we have not measured their \emph{product} at real
sibling-divergence positions in a live rollout. Until that measurement exists,
Table~\ref{tab:anticorr} is a prediction rather than an observation.

\paragraph{Scale.}
All experiments use a single $1.5$B model. The relative size of the two channels
plausibly depends on how sharp the base policy already is, and we have no evidence
about larger models.

\section*{Ethics Statement}

This work studies an optimization-time regularizer for reinforcement learning on
mathematics problems with automatically verifiable answers. It introduces no new
data collection, no human subjects, and no new model release. The intervention acts
on which of several equally-scored continuations receives a slightly smaller
gradient weight; it does not alter the reward, and so does not change what behavior
is being optimized for. The training data is a public competition-mathematics set
and the evaluation set is a public examination.

The main foreseeable risk is indirect. Methods that preserve diversity during RL
fine-tuning are not specific to mathematics, and the same mechanism would apply to a
reward model encoding any objective, including a harmful one; we note that our
method inherits, and does not mitigate, whatever the reward function encodes. The
computational cost of the method (a forward pass through a set of small prediction
heads at every optimization step) is a real environmental cost, which we quantify in
Section~\ref{sec:res:cost} so that it can be weighed against the benefit.

\bibliography{custom}

\appendix

\section{Closed Forms for Section~\ref{sec:blind}}
\label{sec:appendix-closed}

\TODO{Derivation of Eq.~\eqref{eq:secondorder} and the numerical table of first-order
prediction versus actual entropy change for peaked and uniform distributions.}

\section{Implementation Notes}
\label{sec:appendix-impl}

\TODO{Tensor layout of the sibling mask, the one-position alignment shift, the
chunked entropy reduction that avoids materializing the head logits, and the
$\lambda = 0$ bitwise-equivalence test.}

\end{document}
